% concatenated from Goldbach_JNT_Final_Submission.tex
\documentclass[12pt,a4paper]{amsart}
\usepackage{amsfonts}
\usepackage{amsthm}
\usepackage{amsmath}
\usepackage{amssymb}
\usepackage{amscd}
\usepackage[latin2]{inputenc}
\usepackage{t1enc}
\usepackage[mathscr]{eucal}
\usepackage{indentfirst}
\usepackage{graphics}
\usepackage{pict2e}
\usepackage{epic}
\numberwithin{equation}{section}
\usepackage[margin=2.0cm]{geometry}
\usepackage{epstopdf} 
\usepackage{mathtools} 

\usepackage{listliketab}

\usepackage{hyperref}
\hypersetup{
  colorlinks   = true,    
  urlcolor     = black,    % Colour for external hyperlinks
  linkcolor    = black,    % Colour of internal links
  citecolor    = black      % Colour of citations
}









%\documentclass{amsart}
%\documentclass{article}
%\usepackage{amsmath, amsfonts, amssymb, teztcomp}


\theoremstyle{theorem}
\theoremstyle{definition}
\newtheorem{definition}{Definition}
\newtheorem{subdefinition}{Definition}[definition]
\renewcommand{\thesubdefinition}{\thedefinition\alph{subdefinition}}
\newtheorem{theorem}{Theorem}

\newtheorem{case}{Case} 
\newtheorem{cor}{Corollary} 
\newtheorem{example}{Example} 
\newtheorem{proposition}{Proposition} 
\newtheorem{lemma}{Lemma}
\newtheorem{remark}{Remark}
\newtheorem{notation}{Notation}
%\renewcommand\qedsymbol{Q.E.D.}
\newtheorem{corollary}{Corollary}


\theoremstyle{remark}
\theoremstyle{con}




%\newtheorem{remark}[theorem]{\teztbf{Remark}}
%\newtheorem{hyp}{\teztbf{Hypothesis}}
%\newtheorem{question}[theorem]{\teztbf{Question}}
%\newtheorem{notation}[theorem]{\teztbf{Notation}}
\newtheorem{con}[theorem]{Conjecture} 
%\newtheorem{example}[theorem]{Example} 

\newcommand{\Mod}[1]{\ (\mathrm{mod}\ #1)}


\usepackage{nomencl} 


\begin{document}
\title{An Algebraic Proof of the Goldbach Conjecture}
\author{Jason R. South \\
\lowercase{email: jrsouth@smu.edu}
\\\today{\vspace{-5ex}}}

\thispagestyle{plain}
\maketitle
\pagestyle{plain}




%\tableofcontents
\maketitle



% \centerline{\smallit Received: , Revised: , Accepted: , Published: }


\begin{abstract}
We prove the Goldbach Conjecture using p-adic analysis and algebraic methods, requiring no knowledge of prime gaps or distribution. To begin, we define the set of primes up to $a \in \mathbb{N}$ as $\mathcal{P}_a$. It will be shown that if a counter-example $2a$ exists, there exists a polynomial 
\[
 \mathcal{G}_-(z) = \prod_{p_k \in \mathcal{P}_a} (z - p_k) - \prod_{p_k \in \mathcal{P}_a}p_k^{\alpha_k} 
 \]
with conditions $\mathcal{G}_-(2a) = 0$ and all $\alpha_k \in \mathbb{N} \cup \{0\}$. We then use Hensel's Lemma to assist in proving this polynomial structure places the constraint that the solutions for $2a$ require 
\[
S_a = \{p_i \in \mathcal{P}_a : 2a = p_i^{\alpha_i} + p_i \}
\]
which is impossible for $a > 4$ as this requires the primorial $a\#$ to satisfy $a\# \; |\; 2a$.
\end{abstract}


 \noindent\small
\textbf{Keywords:} Goldbach Conjecture, $p$-adic valuation, Hensel's Lemma.
\\
\textbf{MSC Classifications:} 
\textit{Primary}: 11P32, \textit{Secondary}: 11D41, 11S05, 11E95
\normalsize


\pagestyle{headings}

%\markright{\smalltt INTEGERS: 26 (2026) \hfill}
\thispagestyle{empty}
\baselineskip=12.875pt



%\tableofcontents
%%%%%%%%%%%%%%%%%%%%%%%%%%%%%%%%%%%%%%%%%%%%%%%%%%%%%%%








\section{Preliminaries}


\subsection{Necessary Definitions and Auxiliary Theorems}\label{pre}

The Goldbach Conjecture is abbreviated (G.C.) and ":" means "such that" where applicable. Standard conventions apply.

\begin{definition}\label{pset}
The set of primes is $\mathbb{P}$ and $a \in \mathbb{N}_{\geq 2}$. The \textit{Prime Set} of $a$ is $\mathcal{P}_a = \{p_i \in \mathbb{P} : p_i \leq a \}$. The cardinality of this set is given by the prime counting function as $|\mathcal{P} _a| = \pi(a)$.
\end{definition}

\begin{theorem}[\bf\textit{Repeated Root Criterion \cite{dummit2004}}]\label{repeated}
\textit{A root $r$ of $f(x)$ is repeated iff $f'(r) = 0$.}
\end{theorem}

\begin{theorem}[\bf\textit{Hensel's Lemma \cite{Z}}]\label{HL}
Suppose $f(x)$ is a polynomial with integral coefficients. If $f(a) \equiv 0 \Mod{p^j}$ and $f'(a) \not \equiv 0 \Mod{p}$, there is a unique $t \Mod{p} : f(a + tp^j) \equiv 0 \Mod{p^{j + 1}}$.
\end{theorem}



 


\begin{definition}\label{polyx}
A \textit{Goldbach Polynomial} (G.P.) exists when it is possible to take some $a \in \mathbb{N}$ and for each $p_i \in \mathcal{P}_a$ a unique $\alpha_i \in \mathbb{N} \cup \{0\}$ exists producing a mapping $\mathcal{G}_- : \mathbb{C} \to \mathbb{C}$ where  
\begin{equation}\label{roots}
 \mathcal{G}_-(z) = \prod_{p_k \in \mathcal{P}_a} (z - p_k) - \prod_{p_k \in \mathcal{P}_a}p_k^{\alpha_k}
 \end{equation}
 with the condition $\mathcal{G}_-(2a) = 0$. Let the set of roots for the G.P. be given by 
\begin{equation}\label{R}
\mathcal{R}_a = \{ r_i \in \mathbb{C} : \mathcal{G}_-(r_i) = 0 \}
\end{equation}
the factorizations below will be needed in future sections. 
 \begin{equation}\label{factor1}
 \prod_{p_k \in \mathcal{P}_a}(r_i - p_k) = \prod_{p_k \in \mathcal{P}_a}p_k^{\alpha_k} \quad \text{where} \quad -(p_i - r_i) \prod_{p_k \in \mathcal{P}_a \setminus\{p_i\}}(r_i - p_k) = -\prod_{p_k \in \mathcal{P}_a}p_k^{\alpha_k}.
 \end{equation}
\end{definition}



\begin{remark}\label{pderiv}
The Derivative of \ref{roots} follows via the product rule
\begin{equation}
 \mathcal{G}_-'(z) = \prod_{p_k \in \mathcal{P}_a \setminus \{p_1\}} (z - p_k)  +\cdots + \prod_{p_k \in \mathcal{P}_a \setminus \{p_{\pi(a)}\}} (z - p_k)    = \sum_{p_s \in \mathcal{P}_a} \prod_{p_k \in \mathcal{P}_a \setminus \{p_s\}} (z - p_k) 
 \end{equation}
 Note, for any $p_i \in \mathcal{P}_a$ at $z = p_i$ only one term in the equation above remains non-zero, showing
 \begin{equation}\label{roots1}
  \mathcal{G}_-'(p_i) =  \prod_{p_k \in \mathcal{P}_a \setminus \{p_i\}} (p_i - p_k)\quad \text{where} \quad  \mathcal{G}_-'(p_i)  \equiv \mathcal{G}_-'(0)  \prod_{p_k \in \mathcal{P}_a \setminus \{p_i\}} (- p_k)\not \equiv 0 \Mod{p_i}.
 \end{equation} 
Note that since $2a - p_k > 1$ for each $p_k \in \mathcal{P}_a$ showing $\mathcal{G}_-'(2a) > 1$ where Theorem \ref{repeated} shows $2a$ is not a repeated root whenever a G.P. exists.
Therefore, there exists a $Q(z) \in \mathbb{Z}[z]$ where 
\begin{equation}\label{factor}
\mathcal{G}_-(z) = (z - 2a)Q(z) : Q(2a) \neq 0 \quad \text{whose derivative is} \quad \mathcal{G}_-'(z) = Q(z) + (z - 2a)Q'(z).
\end{equation}
\end{remark}




 
\begin{proposition}\label{gbf}
\textit{Let $a \in \mathbb{N}_{> 3}$. If $2a$ is a counterexample to the G.C. then there exists a G.P.}
\end{proposition}

\begin{proof}
Assume $2a$ is counter-example to the G.C. Then no $2a - p_k$ may be prime when $p_k \in \mathcal{P}_a$. If any $2a - p_k$ had a prime factor $a < p < 2a$ then $0< 2a - p_k = n p < 2a$ with $p > a$ forcing $0<n < 2$. Hence, $2a - p_k = p$, contradicting the assumption that $2a$ is a counter-example to the G.C. Therefore, any counter-example requires each $2a - p_k$ only to be divisible by primes in $\mathcal{P}_a$ where the Fundamental Theorem of Arithmetic ensures unique $\alpha_1, \dots, \alpha_{\pi(a)} \in \mathbb{N}\cup \{0\}$ exist satisfying equation \ref{factor1}. The existence of these $\alpha_i$ produce a G.P. from Definition \ref{polyx}.
\end{proof}
 
 
 
 \subsection{A Proof of the Goldbach Conjecture}
  
 We proceed under the assumption that $a \in \mathbb{N}_{>3}$ and fixed where $2a$ is assumed to be a counter-example to the Goldbach Conjecture (G.C.) We then prove that no G.P.s exist beyond $a = 3$ proving the G.C. is true via Proposition \ref{gbf}.
 
 
\begin{definition}\label{pr}
Assume $p_i \in \mathcal{P}_a$. From equation \ref{roots} of Definition \ref{polyx} the G.P. at $z = p_i$ shows $\mathcal{G}_-(p_i) = - \prod_{p_k \in \mathcal{P}_a}p_k^{\alpha_k}$ where the $v_{p_i}(\mathcal{G}_-(p_i)) = \alpha_i$ and $\mathcal{G}_-(p_i) \equiv 0 \Mod{p_i}$ iff $\alpha_i > 0$. Under Remark \ref{pderiv} it was shown $\mathcal{G}_-'(p_i) \not \equiv 0 \Mod{p_i}$. Therefore, via Hensel's Lemma \ref{HL}, if $\alpha_i > 0$,
\begin{equation}\label{pr0}
\mathcal{G}_-(p_i) \equiv 0 \Mod{p_i^{\alpha_i}} \quad \text{and} \quad \mathcal{G}_-(p_i + \tau_{\alpha_i} p_i^{\alpha_i}) \equiv 0 \Mod{p_i^{\alpha_i + 1}}
\end{equation} 
for a unique $p_i$-adic root. This root is the \textit{Principal Root Modulo $p_i$} (P.R.). 
\end{definition}

\begin{lemma}\label{iff}
\textit{$\alpha_i > 0$ if and only if there exists at least one $p_m \in \mathcal{P}_a : 2a \equiv p_m \Mod{p_i}$.}
\end{lemma}

\begin{proof}
Since $2a$ is a root of the G.P. it must satisfy equation \ref{factor1} where $\alpha_i > 0$ if and only if $p_i$ divides at least one $2a - p_m$ term which is true if and only if $2a \equiv p_m \Mod{p_i}$.
\end{proof}


\begin{remark}\label{func}
Definition \ref{pr} showed $\mathcal{G}_-(p_i) = - \prod_{p_k \in \mathcal{P}_a}p_k^{\alpha_k} : p_i \in \mathcal{P}_a$ and using equation \ref{factor} it follows $\mathcal{G}_-(p_i) =  (p_i - 2a)Q(p_i)$. Equation \ref{factor1} then allows, via transitivity, the relation
 \begin{equation}\label{series}
\mathcal{G}_-(p_i) = - \prod_{p_k \in \mathcal{P}_a}p_k^{\alpha_k} = (p_i - 2a)Q(p_i) = (p_i - 2a)\prod_{p_k \in \mathcal{P}_a \setminus \{p_i\}}(2a - p_k). 
\end{equation}
 Note, we may conclude from equation \ref{series} by cancellation of the $p_i - 2a$ term that
 \begin{equation}\label{Q0}
 Q(p_i) = \prod_{p_k \in \mathcal{P}_a \setminus \{p_i\}}(2a - p_k).
 \end{equation}
 Definition \ref{pr} states there exists a unique root satisfying equation \ref{pr0} whenever $\alpha_i > 0$. From Remark \ref{pderiv} the $\mathcal{G}_-(z) = (z-2a)Q(z)$ where equation \ref{pr0} allows 
\begin{equation}\label{series1}
\mathcal{G}_-(p_i + t_{\alpha_i}p_i^{\alpha_i}) \equiv (p_i + t_{\alpha_i}p_i^{\alpha_i}-2a)Q(p_i + t_{\alpha_i}p_i^{\alpha_i}) \equiv 0 \Mod{p_i^{\alpha_i +1}}.
\end{equation}
The question to answer now is that of the relationship between $Q(z)$ and the root $2a$. Note, equation \ref{factor} has $2a$ factored in the G.P., and since $2a$ is not a repeated root, $Q(2a) \neq 0$. Using equations \ref{R}, \ref{factor}, and \ref{Q0} allows for $Q(z)$ to be written in the roots excluding $2a$ where
\begin{equation}\label{Q}
Q(z) = \prod_{r_k \in \mathcal{R}_a \setminus \{2a\}}(z - r_k)^{\sigma_k} \quad \text{where} \quad Q(p_i) = \prod_{r_k \in \mathcal{R}_a \setminus \{2a\}}(p_i - r_k)^{\sigma_k} = \prod_{p_k \in \mathcal{P}_a \setminus \{p_i\}}(2a - p_k).
\end{equation}
This will aid in proving $2a$ is always the P.R. when $\alpha_i > 0$ in the following lemma. We will show in the proof of the following lemma that when $2a$ is assumed not be a P.R. that it must be contained in $Q(z)$, leading to contradiction as this would require $2a$ to be a repeated root. 
\end{remark}

\begin{lemma}\label{perfect}
\textit{For any $p_i \in \mathcal{P}_a$ it follows that if $\alpha_i > 0$, then $2a$ is the Principal Root.}
\end{lemma}

\begin{proof}
Assume $\alpha_i > 0$ and, for contradiction, $v_{p_i}(2a - p_i) = 0$. Therefore, from equation \ref{series1} the  $v_{p_i}(p_i + t_{\alpha_i}p_i^{\alpha_i}-2a) = 0$ where $Q(p_i + t_{\alpha_i}p_i^{\alpha_i}) \equiv 0 \Mod{p_i^{\alpha_i +1}}$ for a unique P.R. in Definition \ref{pr}. The concern is actually not with the construction of the P.R. in this case, but instead with the $p_i$-adic representation of $2a$. Under Remark \ref{pderiv}, $2a$ cannot be a root of $Q(z)$.  However, it is assumed $v_{p_i}(2a - p_i) = 0$ where Lemma \ref{iff} states there exists at least one $p_m \in \mathcal{P}_a \setminus \{p_i\}$ where $2a \equiv p_m \Mod{p_i}$. This requires $Q(p_i) \equiv 0 \Mod{p_i}$ in equation \ref{Q} when $\alpha_i > 0$. Therefore, a unique, finite $p_i$-adic representation exists where $2a = p_m + t_{\mu_i} p_i^{\mu_i} + \cdots$ and $\mu_i \leq \alpha_i$. Furthermore, this $p_i$-adic representation for $2a$ must be in $Q(z)$ from equation \ref{Q}, meaning the root $2a$ is contained in $Q(z)$, contradicting equation \ref{factor}. Therefore, when $\alpha_i > 0$ equation \ref{series1} requires $2a \equiv p_i + t_{\alpha_i}p_i^{\alpha_i} \Mod{p_i}$ to ensure consistency of equation \ref{factor}.
\end{proof}



\begin{theorem}\label{gbt}
\textit{The Goldbach Conjecture is true.}
\end{theorem}

\begin{proof}
Under Lemma \ref{perfect} if $\alpha_i > 0$, then Definition \ref{pr} shows $2a \equiv p_i + t_{\alpha_i}p_i^{\alpha_i} \Mod{p_i^{\alpha_i + 1}}$. Since the $p_i$-adic valuation of $2a - p_i$ is $\alpha_i$ whenever $\alpha_i > 0$, under equation \ref{factor1} no prime may divide multiple $2a - p_k$ when $\alpha_i > 0$. Since all $2a - p_k > 1$ and must be a composition of only primes in $\mathcal{P}_a$, the pigeonhole principal shows the solution set for $2a$ in terms of $\mathcal{P}_a$ is
\[
S_a = \{p_i \in \mathcal{P}_a : 2a = p_i^{\alpha_i} + p_i \}
\]
showing that each $\alpha_i > 0$ and $a\# | 2a$. This is false when $a > 3$ since $2a < a\#$ for $a > 4$, proving no G.P.s exist beyond this value and the Goldbach Conjecture is true via Proposition \ref{gbf}.
\end{proof}



\thispagestyle{plain}


\begin{thebibliography}{1}


\bibitem{Z} Niven, Ivan, Herbert S. Zuckerman, and Hugh L. Montgomery. An Introduction to the Theory of Numbers. 5th ed., John Wiley \& Sons, 1991. p.87,  Theorem 2.23.

\bibitem{dummit2004}
D. S. Dummit and R. M. Foote, \textit{Abstract Algebra}, 3rd ed. Hoboken, NJ: John Wiley \& Sons, Inc., 2004, Section 13.5, Proposition 33.

\end{thebibliography}

\end{document}

%%% Local Variables:
%%% mode: latex
%%% TeX-master: t
%%% End:
